\section*{第 \ref{ch:b1:organic-redox} 章——有机化学中的氧化与还原}

\begin{solution}{exo:b1:organic-redox:1}
乙醇：\ce{CH3} $-3$，\ce{CH2OH} $-1$。乙醛：\ce{CH3} $-3$，CHO $+1$。
乙酸：\ce{CH3} $-3$，COOH $+3$。丙酮：\ce{CH3} $-3$（两个），
C=O $+2$。甲酸：$+2$。
\end{solution}

\begin{solution}{exo:b1:organic-redox:2}
丙-1-醇：先生成丙醛，再生成丙酸。丙-2-醇：生成丙酮。
2-甲基丙-2-醇：不反应（橙色保持不变）。
\end{solution}

\begin{solution}{exo:b1:organic-redox:3}
\ce{CH3COCH3 + 2H+ + 2e- -> CH3CH(OH)CH3};
\ce{CH3CHO + 2H+ + 2e- -> CH3CH2OH}.
\end{solution}

\begin{solution}{exo:b1:organic-redox:4}
\ce{NaBH4}：丁-1-醇、丁-2-醇、不反应、不反应（酸只被
去质子化）。\ce{LiAlH4}：丁-1-醇、丁-2-醇、丁-1-醇和
乙醇、丁-1-醇。
\end{solution}

\begin{solution}{exo:b1:organic-redox:5}
\ce{5CH3CH(OH)CH3 + 2MnO4- + 6H+ -> 5CH3COCH3 + 2Mn^2+ + 8H2O}.
$n = 1.20/60.0 = \qty{0.0200}{mol}$；高锰酸根 $\frac25 \times 0.0200 =
\qty{8.0e-3}{mol}$，需 \qty{0.020}{mol/L} 的溶液 \qty{400}{mL}。
\end{solution}

\begin{solution}{exo:b1:organic-redox:6}
将醇与\omterm{def:b1:nernst:couple}{氧化剂}置于装有蒸馏装置的烧瓶中，在两种沸点之间的温度下加热：
丙醛（\qty{48}{\celsius}）
随其生成即被蒸出，避免进一步氧化；丙-1-醇
（\qty{97}{\celsius}）留在烧瓶中。
\end{solution}

\begin{solution}{exo:b1:organic-redox:7}
\ce{BH4-} 的一对 B--H 键电子进攻羰基碳，C=O 的 $\pi$ 电子对
转移到氧上；\omterm{def:b1:alcohol-activation:alkoxide}{醇盐}从乙醇夺取一个质子。产物中碳上的氢
（\ce{CH3CH2OH}，C1 上的两个氢之一）来自试剂，
氧上的氢来自溶剂。
\end{solution}

\begin{solution}{exo:b1:organic-redox:8}
丁-2-醇具有\omterm{def:b1:stereochemistry-in-depth:stereocentre}{手性}，但平面酮的两个面被进攻的机会均等：
得\omterm{def:b1:stereochemistry-in-depth:enantiomers}{外消旋混合物}，无\omterm{def:b1:stereochemistry-in-depth:optical-activity}{旋光性}。
\end{solution}

\begin{solution}{exo:b1:organic-redox:9}
\ce{3CH3CH2OH + 2Cr2O7^2- + 16H+ -> 3CH3COOH + 4Cr^3+ + 11H2O}：橙色的
重铬酸根（铬为 $+$VI）变为绿色的铬（III）。
\end{solution}

\begin{solution}{exo:b1:organic-redox:10}
2-甲基丁-2-醇：乙醛与 \ce{CH3CH2MgBr} 反应得丁-2-醇（
加成反应，无氧化还原）；氧化得丁酮；丁酮与 \ce{CH3MgBr} 反应，
再经后处理得产物。共一次氧化。戊-3-醇：丙醛与 \ce{CH3CH2MgBr}
一步反应，无氧化还原。
\end{solution}

\begin{solution}{exo:b1:organic-redox:11}
1. 将酮保护为环状\omterm{def:b1:acetals-protection:hemiacetal}{缩醛}（乙烷-1,2-二醇、酸、迪安–斯塔克装置）。
2. 在无水乙醚中用 \ce{LiAlH4} 将酯还原为
伯醇（并生成甲醇）。3. 酸性水溶液脱除\omterm{def:b1:acetals-protection:hemiacetal}{缩醛}：
得 5-羟基戊-2-酮，\ce{CH3CO(CH2)3OH}。
\end{solution}

\begin{solution}{exo:b1:organic-redox:12}
\ce{RCH2OH}：羟基碳的\omterm{def:b1:nernst:oxidation-number}{氧化数}为 $-1$；\ce{RCOOH}：$+3$；共失去四个电子，即
\ce{RCH2OH + H2O -> RCOOH + 4H+ + 4e-}。甲烷（$-4$）氧化为二氧化碳
（$+4$）：共失去八个电子，\ce{CH4 + 2H2O -> CO2 + 8H+ + 8e-}。
\end{solution}

\begin{solution}{pb:b1:organic-redox:1}
\textbf{1.} C1（连接 OH）：0；五个 \ce{CH2}：$-2$。
\textbf{2.} C1（C=O）：$+2$；五个 \ce{CH2}：$-2$。
\textbf{3.} 两个 COOH 碳：$+3$；四个 \ce{CH2}：$-2$。
\textbf{4.} C1（$0 \to +3$）及其邻位的 C6（$-2 \to +3$），后者
成为第二个羧基碳；C1--C6 键断裂。
\textbf{5.} $3 + 5 = 8$ 个电子。
\textbf{6.} 硝酸的氧化能力足以断开与羰基相邻的 C--C 键，
而较温和的\omterm{def:b1:nernst:couple}{氧化剂}做不到。
\textbf{7.} \ce{C6H12O + 3H2O -> C6H10O4 + 8H+ + 8e-}.
\textbf{8.} $+$V 与 $+$I。
\textbf{9.} \ce{2HNO3 + 8H+ + 8e- -> N2O + 5H2O}.
\textbf{10.} 两式相加，八个电子、\ce{8H+} 与三个水分子
相互抵消：\ce{C6H12O + 2HNO3 -> C6H10O4 + N2O + 2H2O}。
\textbf{11.} 8.
\textbf{12.} $10^6/146.0 = \qty{6.85e3}{mol}$ 的酸；$2 \times 6.85 \times
10^3 \times 63.0 = \qty{8.6e5}{g}$，即 \qty{0.86}{t} 硝酸。
\textbf{13.} \ce{H2O2 + 2H+ + 2e- -> 2H2O}.
\textbf{14.} \ce{C6H12O + 4H2O2 -> C6H10O4 + 5H2O}.
\textbf{15.} $4 \times 6.85 \times 10^3 \times 34.0 = \qty{9.3e5}{g}$，
\qty{0.93}{t}.
\textbf{16.} 水：无氮氧化物。
\textbf{17.} 平衡常数很大并不能说明\omterm{def:b1:rate-laws:rate}{反应速率}；
无\omterm{def:b1:elementary-steps:catalyst}{催化剂}时 \ce{H2O2} 反应缓慢，其自身分解亦然。
\textbf{18.} 它将环己酮还原回环己醇：
\ce{4C6H10O + BH4- -> B(OC6H11)4-}，然后
\ce{B(OC6H11)4- + 4H2O -> 4C6H11OH + B(OH)4-}.
\textbf{19.} \qty{6.85e3}{mol}.
\textbf{20.} \qty{6.85e3}{mol} 的 \ce{N2O}（每摩尔酸对应一摩尔）。
\textbf{21.} $6.85 \times 10^3/0.95 \times 100.0 = \qty{7.2e5}{g}$, \qty{0.72}{t}.
\textbf{22.} $6.85 \times 10^3 \times 44.0 = \qty{3.0e5}{g}$: \textbf{\qty{0.30}{t}
每吨己二酸对应的 \ce{N2O}}。
\end{solution}
